\documentclass[11pt,a4paper]{article}

\usepackage[margin=2.5cm]{geometry}
\usepackage{amsmath,amssymb}
\usepackage{graphicx}
\usepackage{float}
\usepackage{booktabs}
\usepackage{hyperref}

\title{\textbf{Mini Project 2026}\\
Modelling a 3D Printer}

\author{
Stefano Natalucci \and
Giorgio Ramundo
}

\date{\today}

\begin{document}

\maketitle

\tableofcontents
\newpage

\section{Physical Modelling}
\subsection{Kinematics}

The CoreXY mechanism couples the motion of both motors to the Cartesian
motion of the print head. The linear displacement
of each belt is related to the corresponding motor rotation through the
motor pulley radius $r_m$.

Using the chosen sign convention, the belt displacements satisfy
\begin{equation}
r_m\theta_l = x+y,
\end{equation}
\begin{equation}
r_m\theta_r = x-y.
\end{equation}

Solving for the Cartesian coordinates gives
\begin{equation}
x = \frac{r_m}{2}\left(\theta_l+\theta_r\right),
\end{equation}
\begin{equation}
y = \frac{r_m}{2}\left(\theta_l-\theta_r\right).
\end{equation}

Differentiating with respect to time yields the print-head velocities:
\begin{equation}
\dot{x} = \frac{r_m}{2}
\left(\dot{\theta}_l+\dot{\theta}_r\right),
\end{equation}
\begin{equation}
\dot{y} = \frac{r_m}{2}
\left(\dot{\theta}_l-\dot{\theta}_r\right).
\end{equation}

\end{document}
